\documentclass[11pt, letterpaper]{article}

\usepackage[margin=1in]{geometry}
\usepackage{titlesec}
\usepackage{enumitem}
\usepackage{hyperref}
\usepackage{xcolor}
\usepackage{fontenc}
\usepackage{parskip}
\usepackage{graphicx}

% ── Colours ──────────────────────────────────────────────────
\definecolor{accent}{HTML}{0EA5E9}
\definecolor{dark}{HTML}{0F172A}
\definecolor{muted}{HTML}{64748B}

\hypersetup{
  colorlinks=true,
  linkcolor=accent,
  urlcolor=accent,
}

\titleformat{\section}{\large\bfseries\color{dark}}{}{0em}{}[\vspace{-0.4em}\color{accent}\rule{\linewidth}{0.4pt}]
\titleformat{\subsection}{\normalsize\bfseries\color{dark}}{}{0em}{}

\pagestyle{empty}

\begin{document}

% ── Title ────────────────────────────────────────────────────
\begin{center}
  {\LARGE\bfseries MasterBuild}\\[6pt]
  {\color{muted}\large Autonomous Market Intelligence for Rapid Enterprise Innovation}
\end{center}

\vspace{0.6em}

% ── What it does ─────────────────────────────────────────────
\section{What MasterBuild Does}

Large organisations are full of smart people with good ideas---but turning a spark of
intuition into something tangible usually means weeks of research, cross-functional
meetings, slide decks, and cautious resource allocation. Most ideas never get that
far; they die in a backlog because nobody can justify the headcount to find out
whether the market actually cares.

MasterBuild changes the math. A user types a single sentence---a product hunch, a
new market angle, an internal tool concept---and an autonomous swarm of five AI
research agents fans out across YouTube, X (Twitter), Reddit, Substack, and Brave
Search simultaneously. Each agent operates its own headless browser, watches real
content, reads real conversations, and pulls back concrete evidence: viral formats,
pricing complaints, willingness-to-pay signals, expert analyses, and raw market
sizing data.

Within minutes those findings are synthesised into a living business plan that
evolves in real time on a 3D dashboard. When the research reaches critical mass
across all four social platforms, MasterBuild automatically produces a
build-ready implementation plan---complete with screens, data models, workflows,
monetisation strategy, and launch steps---and hands the whole thing off to a
no-code builder so the team can have a working MVP prototype in their hands the
same afternoon.

No wasted sprints. No speculative planning. Just evidence-based conviction,
delivered fast enough to keep up with the pace of ideas inside a growing company.

% ── B2B value proposition ────────────────────────────────────
\section{The Enterprise Problem We Solve}

In most companies above a certain size, innovation bottlenecks are not about a
lack of creativity---they are about the \emph{cost of validating} creativity.
Product managers, founders, and internal entrepreneurs all face the same friction:

\begin{itemize}[nosep]
  \item \textbf{Opportunity cost of research.} Assigning even a small team to
        investigate a new idea pulls them off revenue-generating work.
  \item \textbf{Slow feedback loops.} By the time a concept survives internal
        review, the market window may have shifted.
  \item \textbf{Risk aversion by default.} Without hard evidence, leadership
        defaults to saying ``not now''---and good ideas quietly expire.
\end{itemize}

MasterBuild collapses the validation cycle from weeks to minutes. It lets
enterprises treat ideation the way modern engineering treats deploys:
\emph{ship fast, measure immediately, iterate or discard with confidence.}
Employees get to follow their intuition, and the company gets a lightweight
process that surfaces real market demand before a single engineering hour is
spent. The result is a culture where experimentation is cheap, fast, and
data-backed---unlocking new revenue streams that would otherwise never leave
the whiteboard.

% ── How sponsors are used ────────────────────────────────────
\section{Sponsor Technologies in Action}

MasterBuild is built on top of three sponsor platforms, each playing a
critical role in the end-to-end pipeline:

\subsection{InsForge --- Full-Stack Agentic Backend}
InsForge is the shared nervous system of the entire application. The Python
orchestrator writes every discovery, agent observation, business plan revision,
and builder output to InsForge's managed PostgreSQL database through a dedicated
runtime client. The Next.js dashboard reads that same data back through the
InsForge SDK and subscribes to nine real-time channels (missions, agents,
discoveries, logs, signals, agent memory, agent thoughts, business plans, and
builder outputs) so the UI updates live as agents work. InsForge also provides
file storage for agent browser screenshots and authentication for both the
frontend and the Python runtime. In short, InsForge is the only shared state
between the two independent runtimes---without it, the agents and the dashboard
could not communicate at all.

\subsection{MiniMax (M2.7) --- The AI Reasoning Engine}
Every piece of intelligence in MasterBuild is generated by MiniMax's M2.7 model.
It powers the browser-use agents that navigate YouTube, X, Reddit, and Substack;
it summarises raw page content into structured discoveries with keywords; it
periodically rewrites the evolving business strategy; it synthesises
cross-platform findings into a scored business plan; and it produces the final
implementation plan---including screens, data models, workflows, and
monetisation---that gets handed off for prototyping. MiniMax is the sole LLM
behind all of these steps, accessed through a single unified completion
interface.

\subsection{Lovable --- From Plan to Prototype in One Click}
Once MasterBuild's research agents have gathered validated evidence from all
four platforms, the system automatically generates a concise, build-ready prompt
and encodes it into a Lovable launch URL. The dashboard surfaces a
``Build in Lovable'' button that opens Lovable with the prompt pre-loaded, so
the user can go from a raw idea to a functional web-app prototype without
writing a single line of code. This is the last mile of the pipeline: Lovable
turns the AI-validated implementation plan into something a team can actually
click through, test, and iterate on immediately.

\subsection{Other Sponsors}
Going Merry, Eragon, Anyshift, and Eleven Labs are valued sponsors of this
hackathon but are not directly integrated into the current MasterBuild
pipeline. We see clear future opportunities---voice-narrated research briefings
via Eleven Labs, infrastructure monitoring through Anyshift as the platform
scales, cross-org intelligence dashboards in Eragon, and mobile distribution
through Going Merry---but we believe in honest reporting: if a sponsor's
technology is not wired into the product today, we will not claim otherwise.

% ── Architecture at a glance ─────────────────────────────────
\section{How It Works}

\begin{enumerate}[nosep]
  \item \textbf{Mission launch.} The user submits a business idea through the
        Next.js dashboard. A mission record is created in InsForge.
  \item \textbf{Agent swarm.} Five agents launch concurrently, each with its
        own headless Chromium instance, stealth fingerprint, and
        platform-specific extraction logic:
        \begin{itemize}[nosep]
          \item \emph{Echo} (YouTube) --- viral formats, view counts, creator monetisation
          \item \emph{Pulse} (X/Twitter) --- pain points, feature requests, pricing complaints
          \item \emph{Thread} (Reddit) --- willingness-to-pay signals, community size
          \item \emph{Ledger} (Substack) --- expert analysis, competitor breakdowns
          \item \emph{Atlas} (Brave Search) --- market sizing and link curation
        \end{itemize}
  \item \textbf{Continuous synthesis.} Every 25--30 seconds, the orchestrator
        rewrites the shared strategy document and the evolving business plan,
        streaming updates to the dashboard in real time.
  \item \textbf{Convergence.} When validated discoveries exist from all four
        social platforms, MiniMax generates a finalised implementation plan
        and Lovable handoff payload.
  \item \textbf{Prototype.} The user clicks ``Build in Lovable'' and has a
        working MVP scaffold within minutes.
\end{enumerate}

% ── Closing ──────────────────────────────────────────────────
\section{Why This Matters}

The companies that win in the next decade will not be the ones with the biggest
R\&D budgets---they will be the ones that can test the most ideas, the fastest,
with the least waste. MasterBuild gives every team inside an organisation the
power to validate a market opportunity in the time it used to take to schedule
the first meeting about it. That is not incremental improvement. That is a
fundamentally different relationship with innovation.

\vfill
\begin{center}
  \small\color{muted}
  Built for the Agent Master Hackathon, April 2025.
\end{center}

\end{document}
